% Einheit 4 von 4 – Ausklammern (bank/terme/e4.jsonl, zusammenbau v0.1)
%% TODO Zweigzeile Teil 1: was hier gelernt wird (Satz in Schülersprache aus dem Einheitstitel) – Einheit 4
\einheitenkopf[e4]{Einheit 4 von 4 · Ausklammern}
\zweigzeile{ab Kl. 8 (am Gymnasium ab Kl. 7) · keine P10-Aufgabe}
\setcounter{aufgabe}{19}

%% TODO Titel: Ich-kann-Satz fehlt in der Bank (Platzhalter: Kettenname) – Nr. 20
%% TODO Anweisung: ein Satz über der ersten Teilaufgabe fehlt in der Bank – Nr. 20
\begin{aufgabe}{Ausklammern}
%% TODO Darstellung für schwach fehlt in der Bank (terme-e4-k1-s1-v1; Abschnitt „Für schwache Schüler“)
\swz{Klammere den größten gemeinsamen Faktor aus: $4x + 8$}{}
%% TODO Darstellung für schwach fehlt in der Bank (terme-e4-k1-s1-v2; Abschnitt „Für schwache Schüler“)
\swz{Klammere den größten gemeinsamen Faktor aus: $5a + 15$}{}
%% TODO Darstellung für schwach fehlt in der Bank (terme-e4-k1-s1-v3; Abschnitt „Für schwache Schüler“)
\swz{Klammere den größten gemeinsamen Faktor aus: $7b + 14$}{}
%% TODO Darstellung für schwach fehlt in der Bank (terme-e4-k1-s1-v4; Abschnitt „Für schwache Schüler“)
\swz{Klammere den größten gemeinsamen Faktor aus: $3z + 21$}{}
%% TODO Darstellung für schwach fehlt in der Bank (terme-e4-k1-s1-v5; Abschnitt „Für schwache Schüler“)
\swz{Klammere den größten gemeinsamen Faktor aus: $9y + 12$}{}
\swfrage{Was bleibt gleich, was ändert sich?}
%% TODO Darstellung für schwach fehlt in der Bank (terme-e4-k1-s2-v1; Abschnitt „Für schwache Schüler“)
\swz{Klammere den größten gemeinsamen Faktor aus: $x^2 + 5x$}{}
%% TODO Darstellung für schwach fehlt in der Bank (terme-e4-k1-s3-v1; Abschnitt „Für schwache Schüler“)
\swz{Klammere den größten gemeinsamen Faktor aus: $6x^2 + 15x$}{}
%% TODO Darstellung für schwach fehlt in der Bank (terme-e4-k1-s4-v1; Abschnitt „Für schwache Schüler“)
\swz{Klammere den größten gemeinsamen Faktor aus: $6x + 6$}{}
%% TODO Darstellung für schwach fehlt in der Bank (terme-e4-k1-s5-v1; Abschnitt „Für schwache Schüler“)
\swz{Klammere den größten gemeinsamen Faktor aus: $3x^2 + 6x + 15$}{}
%% TODO Darstellung für schwach fehlt in der Bank (terme-e4-k1-s6-v1; Abschnitt „Für schwache Schüler“)
\swz{Klammere aus und mache die Probe: $6x + 30$}{}
%% TODO Darstellung für schwach fehlt in der Bank (terme-e4-k1-s7-v1; Abschnitt „Für schwache Schüler“)
\swz[3]{Finde den Fehler und rechne richtig: \rechnung{5x + 20 &= 5 \cdot (x + 20)}}{}
\swz[3]{Klammere so weit wie möglich aus und mache die Probe: $3a^2 + 6ab + 9a$}{}
\end{aufgabe}

%% TODO Titel: Ich-kann-Satz fehlt in der Bank (Platzhalter: Kettenname) – Nr. 21
%% TODO Anweisung: ein Satz über der ersten Teilaufgabe fehlt in der Bank – Nr. 21
\begin{aufgabe}{Ausklammern – Fehler finden}
\swz[3]{Finde den Fehler und rechne richtig: \rechnung{12x - 8 &= 4 \cdot (3x + 2)}}{}
\end{aufgabe}

\uebersichtskasten{Ausklammern: gemeinsamen Faktor vor die Klammer, Probe durch Ausmultiplizieren. \\ 6x + 9 = 3 · (2x + 3) \quad 4a² + 2a = 2a · (2a + 1)}
